% Work/MAT - 1. rocnik - hodina 39 - rozsirený pracovni list (120 uloh).
% Prvni CTYRI strany pouze zadani. Oddeleni klice: \newpage.
% Zadani a reseni jsou novou vlastni tvorbou na zaklade dosavadniho uciva.
\documentclass[11pt,a4paper]{article}
% Spolecny styl pro ucitelske pripravy MAT (Work/MAT)
% Verze 1.1, 2026-10-23
\NeedsTeXFormat{LaTeX2e}
\ProvidesPackage{priprava}[2026/10/23 v1.1 Sablona ucitelskych priprav MAT]

% Zakladni typografie a cestina
\RequirePackage[T1]{fontenc}
\RequirePackage[utf8]{inputenc}
\RequirePackage[czech]{babel}
\RequirePackage{lmodern}
\RequirePackage{microtype}
\RequirePackage{amsmath,amssymb}
\RequirePackage{xcolor}
\RequirePackage{enumitem}
\RequirePackage[most]{tcolorbox}
\RequirePackage{tabularx,booktabs,array}
\RequirePackage{fancyhdr}
\RequirePackage{needspace}
\RequirePackage{geometry}
\RequirePackage[hidelinks]{hyperref}

\geometry{a4paper,top=17mm,bottom=18mm,left=23mm,right=23mm}
\definecolor{notegray}{gray}{0.39}
\setlength{\parindent}{0pt}
\setlength{\parskip}{2.5pt plus 1pt minus 0.5pt}
\setlist[itemize]{topsep=2pt,itemsep=1pt,parsep=0pt,leftmargin=1.45em}
\setlist[enumerate]{topsep=2pt,itemsep=2pt,parsep=0pt,leftmargin=1.65em}
\setcounter{secnumdepth}{0}
\renewcommand{\arraystretch}{1.12}

% Zahlavi: vlevo rocnik | globalni tema, vpravo cislo: nazev hodiny
% Zapati zustava zcela prazdne.
\pagestyle{fancy}
\fancyhf{}
\renewcommand{\headrulewidth}{0.35pt}
\renewcommand{\footrulewidth}{0pt}
\setlength{\headheight}{14pt}
\newcommand{\nastavPripravu}[4]{%
  \fancyhead[L]{\small #1\enspace\textbar\enspace #2}%
  \fancyhead[R]{\small #3: #4}%
}

% Minimalisticke cernobile boxy. Kazda definice zvlast.
\tcbset{commonlesson/.style={enhanced,breakable=false,colback=white,colframe=black,
  colbacktitle=white,coltitle=black,boxrule=0.45pt,
  arc=0.8mm,outer arc=0.8mm,left=3mm,right=3mm,top=1.6mm,bottom=1.6mm,
  toptitle=1mm,bottomtitle=0.7mm,fonttitle=\bfseries,
  attach boxed title to top left={xshift=2.5mm,yshift=-1.15mm},
  boxed title style={colback=white,colframe=white,boxrule=0pt,left=1mm,right=1mm}}}
\newtcolorbox{zadaniBox}[1][]{commonlesson,title={Zadání},#1}
\newtcolorbox{prikladBox}[1][]{commonlesson,title={Příklad},#1}
\newtcolorbox{definiceBox}[1][]{commonlesson,title={Definice},#1}

% Spolecny box pro zadani a jeho reseni, bez mezery mezi castmi.
% Pouziti:
%   \begin{prikladReseniBox}
%     Zadani ulohy.
%     \tcblower
%     \begin{align*} ... \end{align*}
%   \end{prikladReseniBox}
% Spodni cast ma jemny sedy podklad a delici linku; lze v ni sazet
% vice radku matematiky. Obe casti maji jen jeden spolecny ramecek.
\tcbset{solutionlesson/.style={
  commonlesson,
  enhanced jigsaw,
  breakable,
  colbacklower=black!3,
  segmentation style={solid,draw=black!25,line width=0.35pt},
  fontlower=\small\color{notegray},
  before lower={\noindent\textbf{Řešení:}\enspace},
  leftlower=3mm,rightlower=3mm
}}
\newtcolorbox{prikladReseniBox}[1][]{solutionlesson,title={Příklad},#1}
\newtcolorbox{zadaniReseniBox}[1][]{solutionlesson,title={Zadání},#1}

% Pomucky pro autora pripravy: sedive poznamky; \reseni pro starsi pripravy
\newcommand{\poznamka}[1]{\par\begingroup\small\color{notegray}\itshape
  \textnormal{\textbf{Poznámka pro učitele.}}\enspace #1\par\endgroup}
\newcommand{\reseni}[1]{\par\begingroup\small\color{notegray}
  \textbf{Očekávané řešení:}\enspace #1\par\endgroup}
\newcommand{\blok}[1]{\needspace{5\baselineskip}\subsection*{#1}}
\newcommand{\vyrok}[1]{\uv{#1}}

\usepackage{multicol}
\usepackage{mathtools}
\geometry{top=16mm,bottom=14mm,left=17mm,right=17mm}
\nastavPripravu{1. ročník}{Mocniny a odmocniny}{39}{Velký pracovní list}
\setlength{\columnsep}{7mm}
\setlength{\parskip}{1.5pt}
\setlist[enumerate]{topsep=3pt,itemsep=5pt,parsep=0pt,leftmargin=8mm}
\newcommand{\studenthead}[2]{%
  {\large\bfseries #1}\par\vspace{1mm}
  {\footnotesize #2}\par\vspace{2.4mm}}
\newcommand{\twocoltasks}[2]{%
\begin{minipage}[t]{.49\textwidth}
\begin{enumerate}[label=\textbf{\arabic*.},start=#1]\small
#2
\end{enumerate}
\end{minipage}}
\begin{document}
\noindent\textbf{Jméno:}\enspace\rule{72mm}{.3pt}\hfill\textbf{Třída:}\enspace\rule{21mm}{.3pt}
\vspace{2mm}
\begin{center}{\LARGE\bfseries Mocniny a odmocniny -- 120 úloh}\end{center}
\noindent\small Zapiš postup na vlastní papír. Kde je to nutné, uveď také \textbf{podmínky} původního výrazu.\par
\vspace{3mm}
\studenthead{Strana 1 / 4: Mocniny -- od základů k těžším výrazům}{Úlohy 1--30. Zjednoduš nebo vypočítej; v úlohách s proměnnými určuj podmínky.}
\noindent\begin{minipage}[t]{.49\textwidth}
\begin{enumerate}[label=\textbf{\arabic*.},leftmargin=8mm,itemsep=11pt,topsep=3pt,start=1]
\normalsize
\item $2^4\cdot2^7$
\item $\dfrac{7^9}{7^6}$
\item $\bigl((-3)^4\bigr)^3$
\item $(2^3\cdot5)^2$
\item $4^{-2}$
\item $(-2)^{-3}$
\item $-2^{-4}$
\item $\left(\dfrac34\right)^{-2}$
\item $\dfrac{16^3}{4^4}$
\item $\dfrac{9^2\cdot27}{3^5}$
\item $\dfrac{125^2}{25\cdot5^3}$
\item $\dfrac{(-4)^5}{(-4)^2}$
\item $\dfrac{2^{-3}\cdot2^5}{2^{-1}}$
\item $6^0+(-6)^0$
\item $0^5+1^{51}$
\end{enumerate}
\end{minipage}\hfill
\begin{minipage}[t]{.49\textwidth}
\begin{enumerate}[label=\textbf{\arabic*.},leftmargin=8mm,itemsep=11pt,topsep=3pt,start=16]
\normalsize
\item $2^3+2^4$
\item $(-3^2)^2$
\item $-((-3)^2)$
\item $a^4\cdot a^{-7}$
\item $\dfrac{x^7}{x^{-2}}$
\item $(a^2b^{-3})^{-2}$
\item $\dfrac{x^{-2}y^4}{x^3y^{-1}}$
\item $\dfrac{(-a)^4(-a)^3}{(-a)^2}$
\item $\dfrac{(a^3b^2)^2}{(ab)^3}$
\item $\dfrac{(2a^2b^{-1})^3}{4a^{-1}b^2}$
\item $\left(\dfrac{9x^{-2}y^3}{3xy^{-1}}\right)^{-1}$
\item $\dfrac{(x+y)^5(x+y)^{-2}}{(x+y)^2}$
\item $\left(\dfrac{a^2}{b^3}\right)^{-2}\!\left(\dfrac ab\right)^3$
\item $\dfrac{8^2\cdot6^3}{4^3\cdot12}$
\item $\dfrac{125^2\cdot2^{10}}{16^2\cdot50}\cdot\dfrac{25}{5^3}$
\end{enumerate}
\end{minipage}
\clearpage
\studenthead{Strana 2 / 4: Odmocniny a racionální exponenty}{Úlohy 31--60. Výsledek pokud možno zapiš přesně a usměrni zlomky.}
\noindent\begin{minipage}[t]{.49\textwidth}
\begin{enumerate}[label=\textbf{\arabic*.},leftmargin=8mm,itemsep=11pt,topsep=3pt,start=31]
\normalsize
\item $\sqrt{72}$
\item $\sqrt{200}$
\item $\sqrt{12}+\sqrt{27}$
\item $\sqrt{75}-\sqrt{12}$
\item $2\sqrt{18}-\sqrt{50}$
\item $\sqrt8\cdot\sqrt{18}$
\item $\sqrt6\cdot\sqrt{24}$
\item $\dfrac{\sqrt{50}}{\sqrt2}$
\item $\sqrt{\dfrac3{16}}$
\item $\sqrt{(-7)^2}$
\item $\sqrt{(x-4)^2}$
\item $\sqrt{x^4}$
\item $\sqrt{3x^2}$
\item $\sqrt[3]{-54}$
\item $\sqrt[3]{16}+\sqrt[3]{54}$
\end{enumerate}
\end{minipage}\hfill
\begin{minipage}[t]{.49\textwidth}
\begin{enumerate}[label=\textbf{\arabic*.},leftmargin=8mm,itemsep=11pt,topsep=3pt,start=46]
\normalsize
\item $\sqrt[3]{-8x^3}$
\item $\sqrt[3]{a^6b^3}$
\item $3\sqrt5+2\sqrt5$
\item $\sqrt3+\sqrt5$
\item $\sqrt a+\sqrt a$; $a\ge0$
\item $2\sqrt{12}-\sqrt{27}+\sqrt{75}$
\item $\dfrac1{\sqrt2}$
\item $\dfrac3{2\sqrt3}$
\item $\dfrac1{2+\sqrt3}$
\item $\dfrac1{\sqrt5-\sqrt2}$
\item $(\sqrt6+\sqrt2)(\sqrt6-\sqrt2)$
\item $8^{2/3}$
\item $16^{-3/4}$
\item $81^{1/2}\cdot27^{1/3}$
\item $\dfrac{a^{3/2}}{a^{1/2}}$; $a>0$
\end{enumerate}
\end{minipage}
\clearpage
\studenthead{Strana 3 / 4: Doplň pravidla a část důkazu}{Úlohy 61--90. Doplň prázdná místa; u odmocnin hlídej definiční obor.}
\noindent\begin{minipage}[t]{.49\textwidth}
\begin{enumerate}[label=\textbf{\arabic*.},leftmargin=8mm,itemsep=11pt,topsep=3pt,start=61]
\normalsize
\item $a^m\cdot a^n=a^{\,\square}$
\item $\dfrac{a^m}{a^n}=a^{\,\square}$; $a\ne0$
\item $(a^m)^n=a^{\,\square}$
\item $(ab)^n=\square$
\item $\left(\dfrac ab\right)^n=\square$; $b\ne0$
\item $a^0=\square$; $a\ne0$
\item $a^{-n}=\square$; $a\ne0$
\item $\sqrt[n]{ab}=\square$; pro sudá $n$ požaduj $a,b\ge0$.
\item $\sqrt[n]{a/b}=\square$; pro sudá $n$: $a\ge0,b>0$.
\item $\sqrt{a^2}=\square$
\item $u\sqrt a+v\sqrt a=\square$ ($a\ge0$)
\item $a^{m/n}=\square$ ($a>0$)
\item $a^3a^2=(aaa)(aa)=\square$
\item $\dfrac{a^5}{a^2}=\dfrac{a^2\cdot\square}{a^2}=a^3$
\item Doplň důkaz: $(a^2)^3=a^2a^2a^2=a^{\square}$.
\end{enumerate}
\end{minipage}\hfill
\begin{minipage}[t]{.49\textwidth}
\begin{enumerate}[label=\textbf{\arabic*.},leftmargin=8mm,itemsep=11pt,topsep=3pt,start=76]
\normalsize
\item $(ab)^2=(ab)(ab)=\square$
\item $\dfrac{a^2}{a^5}=\dfrac1{\square}$; $a\ne0$
\item $\sqrt{50}=\sqrt{\square\cdot2}=\square\sqrt2$
\item $2\sqrt3+4\sqrt3=\square\sqrt3$
\item $\dfrac1{\sqrt7}=\dfrac{\sqrt7}{\square}$
\item $\dfrac1{\sqrt5+2}=\dfrac{\sqrt5-2}{\square}$
\item $\sqrt{a^2}=a$ pro $a\,\square\,0$
\item $a^{-2}a^2=\square$ ($a\ne0$)
\item $\sqrt[3]{-27}=\square$
\item $a^{3/2}=(\sqrt a)^{\square}$ ($a>0$)
\item $16^{3/4}=(\square)^3=\square$
\item $2^{-3}=\dfrac1{2^{\square}}=\square$
\item $a^{1/n}=\square$ ($a>0$)
\item $\sqrt{9+16}=\square$, ale $\sqrt9+\sqrt{16}=\square$
\item $\sqrt[3]{8\cdot27}=\square\cdot\square=\square$
\end{enumerate}
\end{minipage}
\clearpage
\studenthead{Strana 4 / 4: Posuď, najdi chybu a použij pravidla}{91--102: pravda/nepravda a zdůvodnění. 103--112: oprav první chybu. 113--120: aplikace a výzvy.}
\noindent\begin{minipage}[t]{.49\textwidth}
\begin{enumerate}[label=\textbf{\arabic*.},leftmargin=8mm,itemsep=11pt,topsep=3pt,start=91]
\normalsize
\item $a^2+a^3=a^5$ pro všechna $a$.
\item $(-2)^4=-(2^4)$.
\item $(ab)^3=a^3b^3$ pro reálná $a,b$.
\item $(a+b)^2=a^2+b^2$ obecně.
\item $\sqrt9=\pm3$.
\item Rovnice $x^2=9$ má řešení $x=\pm3$.
\item $\sqrt{x^2}=x$ pro všechna $x$.
\item $\sqrt[3]{-8}=-2$.
\item $4^{-2}=-16$.
\item $\sqrt8+\sqrt{18}=5\sqrt2$.
\item $\sqrt{\frac ab}=\frac{\sqrt a}{\sqrt b}$ při $a\ge0,b>0$.
\item $a^0=1$ pro všechna $a\in\mathbb R$.
\item Chybný postup: $2^3\cdot2^4=2^{12}$.
\item Chybný postup: $3^6/3^2=3^3$.
\item Chybný postup: $(-5)^2=-25$.
\end{enumerate}
\end{minipage}\hfill
\begin{minipage}[t]{.49\textwidth}
\begin{enumerate}[label=\textbf{\arabic*.},leftmargin=8mm,itemsep=11pt,topsep=3pt,start=106]
\normalsize
\item Chybný postup: $-(-5)^2=25$.
\item Chybný postup: $(a^2)^3=a^5$.
\item Chybný postup: $(a+b)^2=a^2+b^2$.
\item Chybný postup: $\sqrt{9+16}=3+4=7$.
\item Chybný postup: $\sqrt{x^2}=x$ pro $x=-3$.
\item Chybný postup: $\sqrt{12}+\sqrt{27}=\sqrt{39}$.
\item Chybný postup: $\frac1{\sqrt3+1}=\frac{\sqrt3+1}2$.
\item Čtverec má obsah $200\,\mathrm{cm}^2$. Urči stranu.
\item Krychle má objem $216\,\mathrm{cm}^3$. Urči hranu.
\item Rám $0{,}9\times1{,}2\,\mathrm m$: délka úhlopříčky?
\item Rezistor $R=2{,}2\,\mathrm{k}\Omega$, proud $I=0{,}03\,\mathrm A$. Spočti $P=RI^2$.
\item Převeď $0{,}000048\,\mathrm m$ na $\mathrm{\mu m}$.
\item $\dfrac{(a^{-2}b^3)^2}{(ab^{-1})^{-2}}$ ($a,b\ne0$)
\item $\dfrac{(\sqrt{50}+\sqrt8)}{\sqrt2}$
\item $\dfrac{\sqrt3+\sqrt2}{\sqrt3-\sqrt2}-\dfrac{\sqrt3-\sqrt2}{\sqrt3+\sqrt2}$
\end{enumerate}
\end{minipage}

\newpage
\begin{center}{\Large\bfseries VÝSLEDKY A ŘEŠENÍ -- oddělená část}\end{center}
\noindent\textbf{Pro tisk žákům použij pouze strany 1--4.} Tato část začíná explicitním příkazem \texttt{\textbackslash newpage}.\par\vspace{2mm}
\section*{Výsledky 1--30: mocniny}
\begin{enumerate}[label=\textbf{\arabic*.},leftmargin=9mm,itemsep=4pt,start=1]
\small
\item $2^{11}$
\item $7^3$
\item $3^{12}$
\item $2^6\cdot5^2$
\item $\tfrac1{16}$
\item $-\tfrac18$
\item $-\tfrac1{16}$
\item $\tfrac{16}{9}$
\item $16$
\item $9$
\item $5$
\item $-64$
\item $8$
\item $2$
\item $1$
\item $24$
\item $81$
\item $-9$
\item $\tfrac1{a^3}$; $a\ne0$
\item $x^9$; $x\ne0$
\item $\tfrac{b^6}{a^4}$; $a,b\ne0$
\item $\tfrac{y^5}{x^5}$; $x,y\ne0$
\item $-a^5$; $a\ne0$
\item $a^3b$; $a,b\ne0$
\item $\tfrac{2a^7}{b^5}$; $a,b\ne0$
\item $\tfrac{x^3}{3y^4}$; $x,y\ne0$
\item $x+y$; $x+y\ne0$
\item $\tfrac{b^3}{a}$; $a,b\ne0$
\item $18$
\item $250$
\end{enumerate}
\section*{Výsledky 31--60: odmocniny}
\begin{enumerate}[label=\textbf{\arabic*.},leftmargin=9mm,itemsep=4pt,start=31]
\small
\item $6\sqrt2$
\item $10\sqrt2$
\item $5\sqrt3$
\item $3\sqrt3$
\item $\sqrt2$
\item $12$
\item $12$
\item $5$
\item $\tfrac{\sqrt3}4$
\item $7$
\item $|x-4|$
\item $x^2$
\item $|x|\sqrt3$
\item $-3\sqrt[3]2$
\item $5\sqrt[3]2$
\item $-2x$
\item $a^2b$
\item $5\sqrt5$
\item Nelze sloučit na jediný člen.
\item $2\sqrt a$ ($a\ge0$)
\item $6\sqrt3$
\item $\tfrac{\sqrt2}2$
\item $\tfrac{\sqrt3}2$
\item $2-\sqrt3$
\item $\tfrac{\sqrt5+\sqrt2}3$
\item $4$
\item $4$
\item $\tfrac18$
\item $27$
\item $a$ ($a>0$)
\end{enumerate}
\section*{Doplnění 61--90: pravidla a důkazy}
\begin{enumerate}[label=\textbf{\arabic*.},leftmargin=9mm,itemsep=4pt,start=61]
\small
\item $m+n$
\item $m-n$
\item $mn$
\item $a^nb^n$
\item $\dfrac{a^n}{b^n}$
\item $1$
\item $\dfrac1{a^n}$
\item $\sqrt[n]a\sqrt[n]b$
\item $\dfrac{\sqrt[n]a}{\sqrt[n]b}$
\item $|a|$
\item $(u+v)\sqrt a$
\item $\sqrt[n]{a^m}$
\item $5$ a $5$: $3+2$ činitelů.
\item $a^3$ ($a\ne0$)
\item $6$
\item $a^2b^2$
\item $a^3$
\item $25$ a $5$
\item $6$
\item $7$
\item $1$ ($5-4$)
\item $\ge$
\item $1$
\item $-3$
\item $3$
\item $2$ a $8$
\item $3$ a $\tfrac18$
\item $\sqrt[n]a$
\item $5$ a $7$ (protipříklad)
\item $2$, $3$, $6$
\end{enumerate}
\section*{Zdůvodnění 91--120: pravdivost, chyby, aplikace}
\begin{enumerate}[label=\textbf{\arabic*.},leftmargin=9mm,itemsep=4pt,start=91]
\small
\item Ne: pro $a=2$ je $4+8\ne32$.
\item Ne: vlevo $16$, vpravo $-16$.
\item Ano (mocnina součinu).
\item Ne: pro $a=b=1$ je $4\ne2$.
\item Ne: $\sqrt9=3$; $\pm3$ řeší $x^2=9$.
\item Ano: $-3$ i $3$.
\item Ne: pro $x=-2$ je vlevo $2$, vpravo $-2$.
\item Ano: $(-2)^3=-8$.
\item Ne: $4^{-2}=1/16$.
\item Ano: $2\sqrt2+3\sqrt2=5\sqrt2$.
\item Ano, při uvedených podmínkách.
\item Ne: $0^0$ zde nedefinujeme.
\item Sčítáme exponenty: $2^7=128$.
\item Odečítáme: $3^4=81$.
\item Je to sudá mocnina: $25$.
\item Minus je před mocninou: $-25$.
\item Exponenty násobíme: $a^6$.
\item Např. $a=b=1$: $4\ne2$; obecně $a^2+2ab+b^2$.
\item Odmocňujeme až po součtu: $\sqrt{25}=5$.
\item Pro $x=-3$ je $\sqrt{x^2}=3=|x|$.
\item Samostatně zjednoduš: $2\sqrt3+3\sqrt3=5\sqrt3$.
\item Rozšiř sdruženým výrazem: $(\sqrt3-1)/2$.
\item $\sqrt{200}=10\sqrt2\,\mathrm{cm}$.
\item $\sqrt[3]{216}=6\,\mathrm{cm}$.
\item $\sqrt{0{,}9^2+1{,}2^2}=1{,}5\,\mathrm m$.
\item $2200\cdot0{,}03^2=1{,}98\,\mathrm W$.
\item $48\,\mathrm{\mu m}$, protože $1\,\mathrm{\mu m}=10^{-6}\,\mathrm m$.
\item $b^4/a^2$; $a,b\ne0$.
\item $7$.
\item $4\sqrt6$ (jmenovatele usměrni).
\end{enumerate}
\end{document}
